\begin{frame}{수용장 손계산 + 풀링이 스케일업하는 기제}
  \small
  \begin{columns}[T]
    \begin{column}{0.48\textwidth}
      입력 32$\times$32 예:
      \begin{tabular}{@{}llrrl@{}}
        \toprule
        층 & 연산 & $k$ & $s$ & $R_l$ \\
        \midrule
        0 & 입력 & — & — & 1 \\
        1 & Conv & 5 & 1 & 1+4$\cdot$1 = \textbf{5} \\
        2 & MaxPool & 2 & 2 & 5+1$\cdot$1 = \textbf{6} \\
        3 & Conv & 5 & 1 & 6+4$\cdot$2 = \textbf{14} \\
        \bottomrule
      \end{tabular}
      \vskip0.6em
      \textbf{단 세 층 만에}, 마지막 합성곱 뉴런은 원본의
      14$\times$14 영역(32$\times$32의 약 19\%)을 ``본다''.
    \end{column}
    \begin{column}{0.52\textwidth}
      \begin{itemize}
        \item 3$\times$3, s=1 만 계속 쌓으면
          $R_l = 1+2l$ 으로 \textbf{선형}
        \item \textbf{풀링(s=2)을 한 번만} 끼워도 그 뒤의
          \textbf{모든 확장이 2배 스케일업}
        \item 10층만 쌓아도 수십 픽셀, 풀링을 섞으면
          \textbf{이미지 전체}가 한 뉴런의 수용장
      \end{itemize}
      \vskip0.5em
      \begin{block}{짝을 이룬다}
        \kb{합성곱}(패턴 감지)이 ``\textbf{무엇을} 보는가''를
        학습하고, \kb{풀링}(수용장 확대·해상도 축소)이
        ``\textbf{얼마나 넓은 시야}로 보게 하는가''를 결정.
      \end{block}
    \end{column}
  \end{columns}
\end{frame}
